% docs/manual/developer/chapters/03-config-subsystem.tex
% 第 3 章：config 子系统

\chapter{config 子系统}

\modname{config} 子系统是从用户输入到执行引擎之间最复杂的一层。它负责：YAML 解析、schema 校验、reference 名解析、新→旧字典翻译、JSON/NML 迁移。本章按数据流顺序逐个解释，最后给"添加新配置字段"的实战步骤。

\section{子系统鸟瞰}

\modname{config} 包含 6 个公开模块：

\begin{itemize}
  \item \modname{schema}：dataclass 定义的配置 schema（120 行，最薄）
  \item \modname{loader}：YAML → dataclass 加载与校验（361 行）
  \item \modname{adapter}：dataclass → evaluator 期待的 legacy 字典（996 行，最重）
  \item \modname{resolver}：reference 名解析（含 LowRes/MidRes 后缀逻辑，314 行）
  \item \modname{migration}：JSON / NML → YAML 转换（429 行）
  \item \modname{legacy\_processors}：v1.x / v2.0 旧字段兼容层（909 行，逐步缩）
\end{itemize}

数据流：

\begin{center}
\small
YAML \(\to\) loader \(\to\) \texttt{OpenBenchConfig} \(\to\) resolver \(\to\) adapter \(\to\) legacy dict \(\to\) runner.local
\end{center}

migration 是侧路（一次性把旧 config 转成新 YAML）。legacy\_processors 是 loader 内部用的兼容层，处理旧字段名向新字段名的迁移。

\section{schema.py：dataclass 设计}

\file{schema.py} 定义 7 个核心 dataclass。完整字段表见用户卷附录 A（自动生成自本文件）：

\begin{itemize}
  \item \texttt{OpenBenchConfig}：顶层容器，组合下面 6 个子 config + 顶层 \yamlkey{metrics} / \yamlkey{scores} 列表
  \item \texttt{ProjectConfig}：21 字段，分 6 组——identification（3）/ spatial-temporal bounds（3）/ target resolution（4）/ runtime（4）/ groupby analysis（3）/ advanced（4）
  \item \texttt{EvaluationConfig}：仅 \yamlkey{variables}（必填非空 list）
  \item \texttt{ReferenceConfig}：\yamlkey{data_root} + \yamlkey{sources: dict[str, str | list[str]]}（每变量单源或多源）
  \item \texttt{SimulationEntry}：10 字段——\yamlkey{model}/\yamlkey{root_dir}（必填）+ 8 个可覆盖 model profile 的可选字段（\yamlkey{data_type} / \yamlkey{tim_res} / \yamlkey{grid_res} / \yamlkey{data_groupby} / \yamlkey{prefix} / \yamlkey{suffix} / \yamlkey{fulllist} / \yamlkey{variables}）
  \item \texttt{ComparisonConfig}：\yamlkey{enabled} + \yamlkey{items}（图种白名单；省略走小默认）
  \item \texttt{StatisticsConfig}：同结构（但 \yamlkey{statistics.items} 省略时是空集，与 comparison 不同）
\end{itemize}

\subsection{设计原则}

\begin{itemize}
  \item \textbf{所有字段类型注解}：用 \texttt{Optional[X]} 表示可省略
  \item \textbf{默认值显式}：\texttt{field(default\_factory=lambda: [...])} 用于可变默认（避免共享引用 bug）
  \item \textbf{frozen=False}：dataclass 默认可变；adapter 在内部以"读+构造新 dict"方式翻译，不修改 cfg
  \item \textbf{字段顺序}：必填字段在前，可选字段在后（dataclass 顺序约束）
  \item \textbf{命名}：原则上 snake\_case；唯独 \yamlkey{IGBP_groupby} / \yamlkey{PFT_groupby} 用 UPPER\_SNAKE（保留学术界缩写惯例，与论文文本一致；\yamlkey{climate_zone_groupby} 因非缩写故仍小写）
  \item \textbf{Inline 注释 = 文档源}：每个字段后的 \texttt{\# ...} 注释由 \modname{docs.manual.scripts.generate\_config\_schema} 通过 \modname{tokenize} 抽取，生成附录 A 表格的"说明"列。改注释 = 改文档
\end{itemize}

\subsection{命名双轨：years vs syear/eyear}

\openbench{} 内有一个特殊的命名规范，记录在 \file{src/openbench/CONVENTIONS.md}：

\begin{itemize}
  \item \textbf{新 schema}（\file{config/schema.py}）：\yamlkey{years: list[int]} = \texttt{[2004, 2010]}
  \item \textbf{Adapter 输出}（\file{config/adapter.py}）：\yamlkey{syear=2004}, \yamlkey{eyear=2010}
  \item \textbf{Evaluator 运行时}（\file{core/evaluation.py}）：从 GeneralInfoReader 读 \yamlkey{syear}/\yamlkey{eyear}
\end{itemize}

\textbf{原因}：evaluator 是从 v1.x Fortran 迁移过来的 Python，内部依赖 \yamlkey{syear}/\yamlkey{eyear} 这两个具体字段名。完全重命名工作量太大，所以约定：新代码用 \yamlkey{years}（list），adapter 边界做翻译。

\warninline{开新功能时如果需要用到时段信息，从 \texttt{cfg.project.years} 读 list；不要直接读 syear/eyear（除非你写的是 evaluator-facing 代码）。}

\section{loader.py：加载流程}

\file{loader.py} 主入口 \texttt{load\_config(path) -> OpenBenchConfig}。流程：

\begin{enumerate}
  \item 路径检查（\texttt{.yaml} / \texttt{.yml}）
  \item \texttt{yaml.load(f, Loader=\_IncludeLoader)} 解析 + \texttt{!include} 展开
  \item 顶层是 dict（不是 list 或 scalar）
  \item \texttt{\_build\_config(raw)} 构造 dataclass
\end{enumerate}

\subsection{!include 实现}

\file{loader.py:35-79} 定义 \texttt{\_IncludeLoader}（继承 SafeLoader）+ \texttt{\_include\_constructor}：

\begin{itemize}
  \item 路径相对于父 yaml 所在目录
  \item Glob 模式：\texttt{!include sim/*.yaml} 匹配多文件，按文件名排序合并字典
  \item 递归 include：被包含的文件也用 \texttt{\_IncludeLoader} 加载，所以嵌套 \texttt{!include} 也工作
  \item 错误处理：文件不存在或解析失败 → \texttt{ConfigError}
\end{itemize}

\subsection{Backward compat：旧字段迁移}

\file{loader.py:140-186} 处理两类旧字段迁移：

\begin{enumerate}
  \item \yamlkey{options} 块中的字段 → \yamlkey{project}（v2.0 → v3.0）
  \item \yamlkey{comparison.\{tim\_res, grid\_res, timezone, weight\}} → \yamlkey{project}（v2.0 早期 → v3.0）
  \item \yamlkey{options.data\_root} → \yamlkey{reference.data\_root}
\end{enumerate}

每一处迁移都打 \texttt{logger.warning} 提示用户运行 \cli{openbench migrate}。

\subsection{各 \_build\_X 校验函数}

\file{loader.py} 的核心是一组 \_build\_X 函数：

\begin{itemize}
  \item \texttt{\_build\_project}：必填三字段 + years 列表校验 + time\_alignment 枚举校验
  \item \texttt{\_build\_evaluation}：variables 必须非空 list
  \item \texttt{\_build\_reference}：strip data\_root + 剩下的 key 必须是 \texttt{var → str | list[str]} 映射；字符串自动按逗号拆为 list（loader.py:174 那段从 v2.x namelist 直接迁移友好）。schema 字段 \texttt{ReferenceConfig.sources: dict[str, str | list[str]]} 反映了这个多源能力
  \item \texttt{\_build\_simulation}：处理 \_defaults 合并 + 每 entry 必须 model + root\_dir
  \item \texttt{\_build\_comparison} / \texttt{\_build\_statistics}：可选 dict
\end{itemize}

\subsection{simulation \_defaults 合并语义}

\file{loader.py:_build\_simulation 与 \_merge\_variable\_defaults}：

\begin{itemize}
  \item 顶层字段：\texttt{merged = \{**defaults, **entry\}}（浅合并，entry 覆盖 defaults）
  \item \yamlkey{variables} 字段：按变量名 key 深合并（\texttt{\_merge\_variable\_defaults}）
\end{itemize}

\section{resolver.py：reference 解析}

\file{resolver.py} 是 \emph{单一事实源}：所有入口（\cli{check}、adapter、GUI）都调它做 reference 名解析，避免"check 通过但运行时绑定不一致"。

\subsection{多 ref 展开：N ResolvedReference per variable}

\modname{resolve_all_references(cfg)} 接收 \texttt{cfg.reference.sources} 字典——值可能是字符串（单源）或 list[str]（多源）——展开后为\emph{每个 (var, source) 对}产生一个 \texttt{ResolvedReference}。所以多 ref 配置下输出列表长度大于变量数。

\modname{config.adapter.iter_task_sources} 拿到 resolved 列表后再做 (sim × ref) 笛卡尔积，让 \modname{runner.local} 看到完整的 (var, sim, ref) 三元组任务集。

\subsection{ResolvedReference}

每个 \emph{(变量, 源)} 对一个解析结果（来自 \modname{config.resolver:25-40}）：

\begin{minted}{python}
@dataclass
class ResolvedReference:
    var_name: str          # 变量名
    source_name: str       # 用户写的（可能是 base 名）
    resolved_name: str     # 实际命中的 catalog entry 名
    ref_ds: Optional[ReferenceDataset]
    var_map: Optional[VariableMapping]
    status: str            # "ok" / "not_found" / "no_variable" / "ambiguous"
    provenance: str        # "registry" / "fallback"（实例级二态）
    message: str
\end{minted}

\subsection{两种模式}

\begin{itemize}
  \item \textbf{Strict}（\yamlkey{strict_reference: true}）：未解析的 reference → 硬错误，\cli{check} / \cli{run} 退出；细到\emph{字段级}的 LOW provenance 也升级为错
  \item \textbf{Lenient}（默认）：未解析 → warning + 用 fallback 默认值，\cli{run} 可继续
\end{itemize}

\subsection{后缀解析（\_LowRes / \_MidRes / \_HigRes）}

\modname{registry.manager:RESOLUTION_SUFFIXES} 定义了 3 个后缀：\texttt{\_LowRes} / \texttt{\_MidRes} / \texttt{\_HigRes}。\modname{registry.manager.get_reference} 解析逻辑：

\begin{enumerate}
  \item \textbf{完整名（带后缀）}：直接精确匹配，不再变换
  \item \textbf{Base 名 + 提供 sim 分辨率}：从所有变体中按 3 条规则选最优：
  \begin{enumerate}
    \item Time：ref 频率 $\geq$ sim 频率（不能用月数据评 hourly sim）
    \item Space：grid\_res 与 sim 的差最小
    \item 同分时：选频率最低的（避免无谓上采样）
  \end{enumerate}
  \item \textbf{Base 名 + 不提供 sim 分辨率}：返回 None，用 \modname{get_resolution_variants} 列出所有可选
\end{enumerate}

决策追溯写在 \texttt{registry.last\_resolve\_reason}（一句人类可读的原因），\cli{check} 输出会引用。

\subsection{Provenance：两层语义}

\openbench{} 的 provenance 在两个层面：

\paragraph{ResolvedReference.provenance：实例级（2 值）}

\begin{itemize}
  \item \texttt{"registry"}：从 catalog/profile 命中
  \item \texttt{"fallback"}：未命中，用最小默认值（fallback ReferenceDataset）
\end{itemize}

\paragraph{ReferenceDataset.\_provenance：字段级（5 值，分 3 档）}

每个字段（\yamlkey{tim_res}、\yamlkey{grid_res} 等）一个来源标记：

\begin{itemize}
  \item \texttt{"profile"} / \texttt{"existing"} / \texttt{"nc"} → 归 \texttt{PROVENANCE\_HIGH}（来自显式声明、registry 持久化、或 NC 自带元信息）
  \item \texttt{"scan"} → \texttt{PROVENANCE\_MEDIUM}（从目录结构推断）
  \item \texttt{"default"} → \texttt{PROVENANCE\_LOW}（兜底默认值，可能不准）
\end{itemize}

\cli{check} 输出图标：
\begin{itemize}
  \item HIGH：默认 \texttt{✓} 不显式标
  \item MEDIUM：\texttt{\textasciitilde}\hspace{0.2em}（青色，"inferred from directory structure"）
  \item LOW：\texttt{⚠}（黄色，"default — not confirmed from NC or profile"）
  \item LOW + \yamlkey{strict_reference}：\texttt{✗}（红色，"unconfirmed default"，触发 SystemExit）
\end{itemize}

\section{adapter.py：新→旧翻译}

\file{adapter.py} 是 \modname{config} 中最大的模块（约 1040 行），但概念简单：把 \texttt{OpenBenchConfig} dataclass 翻译成 evaluator 期待的 legacy 字典格式。

\subsection{为什么需要 adapter}

evaluator 是从 v1.x Fortran 迁移到 Python 的代码（住在 \modname{core/evaluation.py}）。它期待的 input 是嵌套字典：

\begin{minted}{python}
{
  "general": {"basename": "...", "syear": 2010, "eyear": 2014, ...},
  "evaluation_items": {"Evapotranspiration": True, ...},
  "metrics": [...],
  "scores": [...],
  "comparisons": [...],
  "statistics": [...],
  ...
}
\end{minted}

而新 schema 是 dataclass。adapter 做的就是从 dataclass 提取信息生成这种字典。

\subsection{核心入口}

\begin{minted}{python}
from openbench.config.adapter import build_runner_bindings

bindings = build_runner_bindings(cfg)  # cfg: OpenBenchConfig
# bindings.runner_cfg = RunnerConfig(...)         -- 主结构
# bindings.namelists  = LegacyNamelists(...)      -- 翻译后的 main/ref/sim 字典
# bindings.figures    = LegacyFigureConfig(...)   -- 可视化 namelist
\end{minted}

\subsection{关键数据结构}

\begin{itemize}
  \item \texttt{RunnerConfig}（frozen dataclass）：basename / basedir / general / evaluation\_items / metrics / scores / comparisons / statistics
  \item \texttt{BridgeRuntimeInfo}：每个 task 的 reference/simulation 信息
  \item \texttt{NamelistBundle}：翻译后的 main\_nl / ref\_nml / sim\_nml 字典
\end{itemize}

\subsection{LEGACY\_GENERAL\_KEYS}

\file{adapter.py:18-50} 列出 evaluator 期待的所有 legacy general 字段名：basename, basedir, syear, eyear, num\_cores, time\_alignment, IGBP\_groupby 等。每次 evaluator 加新需求字段，这里要同步加。

\warninline{adapter.py 是 \openbench{} 演进中变动最频繁的模块之一。改之前先看 \file{tests/test\_config/test\_adapter.py} 与 \file{test\_adapter\_namelists.py}，每次改完跑这两个测试。}

\section{migration.py：JSON/NML → YAML}

\file{migration.py} 的入口：

\begin{minted}{python}
def migrate_config(main_config_path, output_path) -> dict[str, Any]:
    """读旧主配置 → 解析 ref/sim/var 子配置 → 合并为新 YAML"""
\end{minted}

\subsection{支持的输入格式}

\begin{itemize}
  \item \texttt{.json}：v2.0 主配置（顶层是 dict）
  \item \texttt{.yaml} / \texttt{.yml}：v2.0 后期 YAML
  \item \texttt{.nml}：v1.x Fortran namelist（最早格式）
\end{itemize}

\subsection{多文件解析}

旧格式通常有 4-6 个文件，主配置 reference 子配置路径：

\begin{minted}{text}
main.nml         (主)
  ├─ ref_nml.nml   (reference 详细)
  ├─ sim_nml.nml   (simulation 详细)
  └─ var_def.nml   (变量定义)
\end{minted}

\texttt{migrate\_config} 跟踪所有引用，递归读取，最后合并为单一 YAML。

\subsection{字段重命名表}

\file{migration.py} 内部维护 v1.x/v2.0 → v3.0 的字段名映射。新加的兼容字段在 \modname{legacy\_processors} 里。

\section{legacy\_processors.py：兼容层}

\file{legacy\_processors.py} 是 v1.x/v2.0 字段名的"垃圾回收站"：

\begin{itemize}
  \item v1.x 老字段名（\texttt{Compare\_Time\_Resolution}）→ v3.0（\texttt{tim\_res}）
  \item v2.0 中间名（\texttt{compare\_tim\_res}）→ v3.0（\texttt{tim\_res}）
  \item 已废弃但仍可读的 \texttt{options.\textit{X}} → \texttt{project.\textit{X}}
\end{itemize}

\noteinline{\file{legacy\_processors.py} 909 行；目标是逐版本缩。新功能不要写到这里。修 v3.0 字段时不要触碰旧字段处理（除非用户报告 v2.0 配置 migration 失败）。}

\section{添加新配置字段（5 步实战）}

假设你要加 \yamlkey{project.cache\_strategy: "zstd" | "none"}：

\subsection{Step 1：改 schema.py}

\begin{minted}{python}
@dataclass
class ProjectConfig:
    # ... 现有字段
    cache_strategy: str = "zstd"   # 加默认值
\end{minted}

\subsection{Step 2：改 loader.py}

在 \texttt{\_build\_project} 里加：

\begin{minted}{python}
cache_strategy = raw.get("cache_strategy", "zstd")
if cache_strategy not in ("zstd", "none"):
    raise ConfigError(
        f"project.cache_strategy must be 'zstd' or 'none', "
        f"got '{cache_strategy}'"
    )
return ProjectConfig(
    # ... 其他字段
    cache_strategy=cache_strategy,
)
\end{minted}

\subsection{Step 3：改 adapter.py（如果 evaluator 需要）}

如果新字段要传给 evaluator：

\begin{minted}{python}
# 在 build_runner_bindings 中
general["cache_strategy"] = cfg.project.cache_strategy

# 同时把字段加到 LEGACY_GENERAL_KEYS
LEGACY_GENERAL_KEYS = {
    # ... 现有
    "cache_strategy",
}
\end{minted}

\subsection{Step 4：加测试}

\begin{minted}{python}
# tests/test_config/test_loader.py
def test_cache_strategy_default():
    cfg = load_config(FIXTURES / "minimal.yaml")
    assert cfg.project.cache_strategy == "zstd"

def test_cache_strategy_invalid():
    yaml = "..."  # 含 cache_strategy: bogus
    with pytest.raises(ConfigError, match="cache_strategy"):
        load_config_from_str(yaml)
\end{minted}

\subsection{Step 5：改文档}

\begin{itemize}
  \item 用户卷第 3 章 \yamlkey{project} 字段表中加新字段
  \item 跑 \cli{cd docs/manual \&\& make generated} 让附录 A 自动包括新字段
\end{itemize}

\section{下一步}

下一章深入 \modname{data} 子系统：pipeline 阶段、三层缓存、重网格、单位转换、站点匹配。
